\begin{proof}
Throughout, Theorem \claimref{synthesis.hadamard-dirac-states-scale-to-the-flat-massless-vacuum-at-every-point-of-a-gravitational-field} is cited as (H): its clause (i) gives the form $U/\sigma + V\log\sigma$ with $U = \Delta^{1/2}R$ and $D_{x}\sigma = \hb^{-1}(\dot\gamma(1))$, the physical tangent at $x$ of the geodesic $\gamma$ from $x'$ to $x$, and its clause (ii) the Taylor data. The form of the parametrix is Fact \claimref{synthesis.imported-hadamard-parametrix-of-the-squared-dirac-operator-leads-with-van-vleck-root-times-spin-transport}.

\emph{(i).} $D_{x}(U/\sigma) = (D_{x}\sigma^{-1})U + \sigma^{-1}D_{x}U = -sU/\sigma^{2} + (D_{x}U)/\sigma$, so the leading part of $G(x, x')$ is $-\Delta^{1/2}sR/\sigma^{2}$. Exchanging the points reverses the geodesic: $\sigma$ and $\Delta$ are symmetric, $R(x', x) = \rev{R}$, and $D_{x'}\sigma$ is the physical tangent at $x'$ of the reversed geodesic. The physical tangent is transported by the rotor, $\hb^{-1}(\dot\gamma(1)) = R\,\hb^{-1}(\dot\gamma(0))\,\rev{R}$ (the geodesic equation and $\dot R = -\tfrac12\Omega(\dot\gamma)R$), so $s' = -\rev{R}sR$ and $Q(s') = Q(s) = 2\sigma$. Then
\[
\langle a\,(sR)\,b\,(s'\rev{R})\rangle_{0} = -\langle a\,s\,R\,b\,\rev{R}\,s\,R\,\rev{R}\rangle_{0} = -\langle a\,s\,\hat b\,s\rangle_{0} = -\bigl(2(a\cdot s)(\hat b\cdot s) - (a\cdot\hat b)Q(s)\bigr),
\]
by $R\rev{R} = 1$ and the trace $\langle a\,x\,b\,x\rangle_{0} = 2(a\cdot x)(b\cdot x) - (a\cdot b)Q(x)$ of clause (i) of Theorem \claimref{synthesis.one-loop-vacuum-polarisation-extends-without-a-regulator-and-ward-preserving-extensions-differ-only-by-the-renormalisation-constants}. Multiplying by $\Delta^{1/2}\Delta^{1/2}/\sigma^{4}$ gives $\Pi_{6}$. In $\langle a\,(sR)\,b\,\rev{R}\rangle_{0}$ and $\langle a\,R\,b\,(s'\rev{R})\rangle_{0}$ the argument is a product of three vectors and rotors, an odd element, with no scalar part; and $\langle a\,R\,b\,\rev{R}\rangle_{0} = a\cdot\hat b$. The check verifies these identities with four exact rotors (Cayley transforms of rational bivectors, simple and not), six vectors $s$ and sixteen pairs $a$, $b$, and that the identity fails if the tangent at $x'$ is not transported.

\emph{(ii).} The statement $a\cdot D_{x}R = -\tfrac14\mathcal{R}(a\wedge\xi)R + O(\xi^{2})$ is covariant under rotation gauge transformations ($R \mapsto S(x)R\rev{S}(x')$, $\mathcal{R}(B) \mapsto S\mathcal{R}(\rev{S}BS)\rev{S}$), so it may be verified in one gauge. Take $x_{0}$ with $\hb(x_{0}) = 1$, first derivatives of $\hb$ zero at $x_{0}$ (normal position gauge) and $\Omega(a)(x_{0} + z) = \tfrac12\mathcal{R}(z\wedge a) + O(z^{2})$; this is consistent, since then $a\cdot\nabla\Omega(b) - b\cdot\nabla\Omega(a) = \tfrac12\mathcal{R}(a\wedge b) - \tfrac12\mathcal{R}(b\wedge a) = \mathcal{R}(a\wedge b)$ at $x_{0}$, and any $\Omega$ vanishing at $x_{0}$ differs from it at first order by the symmetric part of $a\cdot\nabla\Omega(b)$, which a rotation gauge transformation quadratic in $z$ removes. By (H)(ii), with $x = x_{0} + ly$, $x' = x_{0} + ly'$, $d = y - y'$, $\bar y = \tfrac12(y + y')$: $R = 1 - \tfrac{l^{2}}{2}(\bar y\cdot\nabla)\Omega(d) + O(l^{3}) = 1 - \tfrac{l^{2}}{4}\mathcal{R}(\bar y\wedge d) + O(l^{3})$. The scaled derivative $l\,a\cdot D = a\cdot\partial_{y} + \tfrac{l}{2}\Omega(a)(x_{0} + ly) + O(l^{2})$ applied to it gives, at order $l^{2}$,
\[
-\tfrac14\bigl[\tfrac12\mathcal{R}(a\wedge d) + \mathcal{R}(\bar y\wedge a)\bigr] + \tfrac14\mathcal{R}(y\wedge a) = -\tfrac14\bigl[\tfrac12\mathcal{R}(a\wedge d) - \tfrac12\mathcal{R}(d\wedge a)\bigr] = -\tfrac14\mathcal{R}(a\wedge d),
\]
using $y - \bar y = \tfrac12 d$; with $\xi = ld$ this is the claim, and it does not depend on $\bar y$. Contracting, $D_{x}R = \gamma^{\mu}(e_{\mu}\cdot D_{x}R) = -\tfrac14\gamma^{\mu}\mathcal{R}(e_{\mu}\wedge\xi)R$, and $\gamma^{\mu}\mathcal{R}(e_{\mu}\wedge\xi) = \partial_{a}\cdot\mathcal{R}(a\wedge\xi) + \partial_{a}\wedge\mathcal{R}(a\wedge\xi) = \mathcal{R}(\xi)$ by the protraction identity. A field strength of Weyl type has $\partial_{a}\mathcal{R}(a\wedge\xi) = 0$ for every $\xi$. By (H)(ii), $\Delta = 1 + \tfrac16\Ric^{\circ}(\xi, \xi) + O(\xi^{3})$ with $\Ric^{\circ}(\xi, \xi) = -\xi\cdot\mathcal{R}(\xi)$, so $\Delta^{1/2} = 1 - \tfrac1{12}\xi\cdot\mathcal{R}(\xi)$ and its gradient is $-\tfrac16\mathcal{R}(\xi)$ (the Ricci map is symmetric). Hence $D_{x}U = (-\tfrac16 - \tfrac14)\mathcal{R}(\xi)R + O(\xi^{2}) = -\tfrac5{12}\mathcal{R}(\xi)R + O(\xi^{2})$. The corrections from $\hb - 1 = O(z^{2})$ and from the deviation of the geodesic from the straight line are of relative order $l^{2}$ and enter $D_{x}R$ at order $\xi^{2}$. The check verifies the derivative formula for a general linear $\mathcal{R}$ at exact points (the symmetric difference of a quadratic polynomial is its derivative), the consistency of the gauge, the contraction on field strengths of type D ($B \mapsto c(B + 3\sigma_{r}B\sigma_{r})$ with $\sigma_{r}^{2} = 1$, the form of the stored Schwarzschild field strength), on sums of rotated copies and on field strengths $a\wedge b \mapsto \tfrac12(P(a)\wedge b + a\wedge P(b))$ of Ricci type, and that a map violating the protraction identity leaves a trivector.

\emph{(iii).} Write $G(x, x') = -\Delta^{1/2}sR/\sigma^{2} + (D_{x}U)/\sigma + mU/\sigma + Vs/\sigma + (D_{x}V + mV)\log\sigma$. By (SD), with $s = O(\xi)$ smooth and vanishing to first order: the first term has degree $3$; $mU/\sigma$ degree $2$; $(D_{x}U)/\sigma$ degree at most $1$ by (ii), and at most $0$ where the Ricci vector map vanishes; $Vs/\sigma$ degree at most $1$, with coefficient $V(x', x')$ plus terms of degree at most $0$; the logarithmic term degree at most $0$. The same holds for $G(x', x)$. In the product: degree $6$ is $\Pi_{6}$; degree $5$ could only come from the leading part times the mass term, which vanishes by (i); degree at most $4$ comes from mass times mass ($m^{2}\Delta\,a\cdot\hat b/\sigma^{2}$), leading times $(D U)/\sigma$ (coefficient the Ricci vector map by (ii)) and leading times $Vs/\sigma$ (coefficient $V(x', x')$, a scalar because the zero-order term of the squared operator is a scalar by the Lichnerowicz formula of the cited Fact, clause (b)); everything else has degree at most $3$. The Weyl field strength is absent from $\Pi_{6}$ by (i), and from the coefficients of the terms of degree $4$ by (ii) and because $\Delta - 1$ is Ricci at second order by (H)(ii). The terms with the smooth part of a Hadamard state have degree at most $3$.

\emph{(iv).} By (i), for fixed $x'$, in the normal coordinates $\xi'$ of $x'$ ($s' = D_{x'}\sigma$ is minus the position vector of $x$ in them, and $Q(s') = 2\sigma$) and with $a$ transported to $x'$, $\Pi^{0} = -(2(\hat a\cdot s')(b\cdot s') - (\hat a\cdot b)Q(s'))/\sigma^{4}$ is the flat massless kernel, a homogeneous function of degree $-6$ of $\xi'$. Fact \claimref{synthesis.imported-epstein-glaser-extension-theorem-and-ward-preserving-time-ordered-products}(ii), which holds locally on a gravitational field with the scaling degree read in local coordinates, then gives: extensions of $\Pi$ of degree $6$ exist and two differ by $\sum_{|\alpha| \le 2}c_{\alpha}\partial^{\alpha}\delta$, $\omega = 6 - 4 = 2$; the stored flat extension of clause (ii) of the cited vacuum-polarisation theorem, read in $\xi'$, extends $\Pi^{0}$ with the stored logarithmic coefficient; the part of degree $4$ has $\omega = 0$; the part of degree at most $3 < 4$ extends uniquely. The number of multi-indices of order at most $2$ in four variables is $1 + 4 + 10 = 15$. Since the Weyl field strength enters only the last part, it adds no freedom.
\end{proof}
